\uuid{8MKY}
\exo7id{4336}
\titre{exo7 4336}
\auteur{quercia}
\organisation{exo7}
\datecreate{2010-03-12}
\isIndication{false}
\isCorrection{true}
\chapitre{Intégration}
\sousChapitre{Intégrale de Riemann dépendant d'un paramètre}
\cpge{mpsi-mp2i.limites-continuite.b, pcsi.limites-continuite.b, ptsi.limites-continuite.b, mpsi-mp2i.analyse-asymptotique.a, pcsi.analyse-asymptotique.a, ptsi.analyse-asymptotique.a, mpsi-mp2i.integration.c, pcsi.integration.b, ptsi.integration.b, bcpst1.limites-continuite.b, bcpst1.integration-segment.a, bcpst1.integration-segment.b}
\module{Analyse}
\niveau{L2}
\difficulte{}

\contenu{
\question{
Soit $f : {[0,+\infty[} \to {\R^+}$ continue telle que
$f(x) \int_{t=0}^x f^2(t)\,d t \to \ell\ne 0$ (lorsque $x\to+\infty$).
Trouver un équivalent de $f$ en~$+\infty$.
}
\reponse{
$\Phi(x) =  \int_{t=0}^x f^2(t)\,d t  \Rightarrow  \Phi'\Phi^2 \to \ell^2
          \Rightarrow  \Phi^3 \sim 3\ell^2 x  \Rightarrow  f = \sqrt{\Phi'} \sim \sqrt[3]{\frac \ell{3x}}$.
}
}
